\documentclass{template}

\usepackage{xeCJK}
\setCJKmainfont[AutoFakeBold=2.5]{IPAexGothic}
\setCJKsansfont[AutoFakeBold=2.5]{IPAexGothic}

\usepackage[square,numbers]{natbib}

\usepackage{hyperref,xcolor}
\hypersetup{
    pdftitle={履歴書 - Alexander Weichart},
    pdfauthor={Alexander Weichart},
    pdfsubject={フルスタックソフトウェアエンジニア - AIネイティブ開発 - 履歴書},
    pdfkeywords={Alexander Weichart, 履歴書, ソフトウェアエンジニア, フルスタック, Java, TypeScript, Spring Boot, React, PostgreSQL, Kubernetes, GitLab, Flux CD, Linux, DevOps, Agile, AI, LLM, Claude Code, Codex, Copilot, MCP, エージェント開発},
    pdfcopyright={Copyright (C) 2026 Alexander Weichart},
    pdfmetalang={ja},
    linkcolor=blue,
    citecolor=blue
}

\setname{Alexander Weichart}{}
\setposition{フルスタックソフトウェアエンジニア | AIネイティブ開発}
\setaddress{大阪府、日本}
\setbirthday{アレクサンダー・ヴァイヒャート}
\renewcommand{\getPhone}{\ifdefempty{\PHONE}{}{電話: \PHONE}}
\setmobile{\getPhone}
\setmail{alexanderweichart@fastmail.com}
\setcitizenship{ドイツ}
\setlinkedinaccount{https://www.linkedin.com/in/alexander-weichart-a357241a4/}
\setprojects{https://alexanderweichart.de/4_Projects/Projects}
\sethomepage{https://www.alexanderweichart.de/}
\setgithubaccount{https://github.com/AlexW00}
\setprofilepic{profile.jpeg}
\setthemecolor{Black}


\begin{document}

\vspace*{3mm}
\headerview
\vspace{3ex}

%
\section{スキル}
    \newcommand{\skillone}{\createskill{AIツール}{Claude Code \cpshalf Codex \cpshalf GitHub Copilot}}
    \newcommand{\skilltwo}{\createskill{エージェント開発}{Skills, sub-agents, MCP \cpshalf コンテキスト設計、モデル選択}}
    \newcommand{\skillthree}{\createskill{主要言語}{Java \cpshalf TypeScript \cpshalf Swift}}
    \newcommand{\skillfour}{\createskill{主要技術}{Spring Boot, React, PostgreSQL \cpshalf Kubernetes, GitLab, Flux \cpshalf Linux}}
    \newcommand{\skillfive}{\createskill{言語能力}{ドイツ語（母語）\cpshalf 英語（流暢）\cpshalf 日本語（日常会話）\cpshalf フランス語（基礎）}}
    %
    \createskills{\skillone, \skilltwo, \skillthree, \skillfour, \skillfive}
    \vspace{-3mm} % reduces whitespace!
    %
\bigskip

    \section{職務経歴}
    %
    \datedexperience{Birchill（大阪、日本）}{2026年4月 - 現在}
    \explanation{フルスタックソフトウェアエンジニア（ワーキングホリデー）- TypeScript, AWS}{大阪、日本}
    \explanationdetail{
    \smallskip
    \coloredbullet\ %
    フルスタックの日本語学習プラットフォーム開発を主導し、フロントエンド、バックエンド、AWSインフラの技術設計からリリースまで担当。

    \smallskip
    \coloredbullet\ %
    Lambda、DynamoDB、Step Functions、CDKを用いたTypeScriptサービスとデータ処理基盤を構築し、辞書データや学習機能を実装。
    }
    %
    \datedexperience{DEBTVISION GmbH（バーデン＝ヴュルテンベルク州立銀行）}{2023年10月 - 2026年3月}
    \explanation{フルスタックソフトウェアエンジニア - Java, TypeScript, Kubernetes}{シュトゥットガルト、ドイツ}
    \explanationdetail{
    \smallskip
    \coloredbullet\ %
    B2B金融プラットフォームにおいて、API設計、データベースモデリング、CI/CDパイプライン構築、Kubernetes本番環境デプロイまで機能をエンドツーエンドで担当。

    \smallskip
    \coloredbullet\ %
    部門横断型AI推進チームの開発代表として、チームで初めてGitHub Copilotを導入。ワークショップとメンタリングを通じてエージェント型開発の定着を主導し、6ヶ月間で\mbox{スプリントベロシティ}80\%向上に貢献。

    \smallskip
    \coloredbullet\ %
    クラウドインフラの最適化により、CPUを2倍・メモリを12\%増強しつつ、コストを17\%削減。

    \smallskip
    \coloredbullet\ %
    本番クラスタを別のクラウドプロバイダーにリカバリクラスタとしてゼロから復元する自動化ツールを構築。
    }
    %
    \datedexperience{レーゲンスブルク大学}{2020年9月 - 2023年9月}
    \explanation{学生リサーチアシスタント - JavaScript, Kotlin, C, Neo4j, Figma}{レーゲンスブルク、ドイツ}
    \explanationdetail{
    \smallskip
    \coloredbullet\ %
    カリキュラムをグラフ構造で整理するプラットフォーム\href{https://hci.ur.de/projects/graphit}{GraphIT}の開発をリードし、DELFI 2025で発表。
    }

\bigskip

\section{\texorpdfstring{プロジェクト\ \projectslinkicon}{プロジェクト}}
    %
    \datedexperience{Weichart Apps\ \href{https://apps.weichart.de}{{\normalfont\mdseries\upshape\textcolor{NavyBlue}{\raisebox{0.15ex}{\faLink}}}}}{}
    \explanationdetail{
    \smallskip
    \coloredbullet\ %
    \href{https://apps.weichart.de}{丁寧に設計したiOSアプリ3本}を開発・公開。紙のような使い心地のカレンダー\href{https://apps.apple.com/app/id6747954401}{Hibi}を含め、累計1.5万回以上ダウンロード。エージェント型開発ワークフローを活用。
    }
    %
    \datedexperience{AIネイティブ・ナレッジシステム\ \href{https://alexanderweichart.de/4_Projects/obsidian-setup-2026/My-AI-native-Obsidian-Setup}{{\normalfont\mdseries\upshape\textcolor{NavyBlue}{\raisebox{0.15ex}{\faLink}}}}}{}
    \explanationdetail{
    \smallskip
    \coloredbullet\ %
    定期実行するClaude Codeルーティン、再利用可能なSkills、MCP/CLI連携、永続状態、承認フローを備えたGitベースのナレッジシステムを設計・構築。公開サイトは月間3,000アクセス超。
    }
    %
    \datedexperience{Obsidian 3D Graph\ \href{https://github.com/AlexW00/obsidian-3d-graph}{{\normalfont\mdseries\upshape\textcolor{NavyBlue}{\raisebox{0.15ex}{\faLink}}}}}{}
    \explanationdetail{
    \smallskip
    \coloredbullet\ %
    ObsidianのノートをWebGLグラフで可視化するプラグイン。ダウンロード数5万件超、GitHubスター350件超。
    }
    %
    \datedexperience{AIプロトタイプ（2023年〜）\ \href{https://alexanderweichart.de/4_Projects/Projects}{{\normalfont\mdseries\upshape\textcolor{NavyBlue}{\raisebox{0.15ex}{\faLink}}}}}{}
    \explanationdetail{
    \smallskip
    \raggedright
    \coloredbullet\ %
    2023年からLLMを活用：自然言語をJavaScriptに変換する\href{https://alexanderweichart.de/4_Projects/Projects\#-gpt-lang-2023-03-24}{GPT-Lang}、\href{https://alexanderweichart.de/4_Projects/Projects\#-tandem-gpt-2023-02-04}{Tandem GPT}、\mbox{最新作の\href{https://alexanderweichart.de/4_Projects/Projects\#ai-hud-for-japanese-text-2025-12-22}{日本語テキスト認識AI-HUD}}。
    }

\bigskip

\section{学歴}
    \datedexperience{文学士 - メディア情報学（GPA: 3.5 / 4.0 米国式換算）}{2019年9月 - 2023年9月}
    \explanation{レーゲンスブルク大学}{レーゲンスブルク、ドイツ}
    \explanationdetail{
        \smallskip
        \coloredbullet\ %
        フルスタック開発、アジャイルプロジェクト管理、自然言語処理、倫理学。
     }

\bigskip

\section{コミュニティ}
\indent
\explanationdetail{
\smallskip
\coloredbullet\ %
複数の人気オープンソースプロジェクトのメンテナー（\href{https://github.com/AlexW00}{GitHub}）。合計スター数1,000件超。

\smallskip
\coloredbullet\ %
日独協会レーゲンスブルク（\href{https://djg-regensburg.de/de/}{DJG Regensburg}）の支援会員。

}
\bigskip

\section{論文・発表}
\indent
\explanationdetail{
\smallskip
\coloredbullet\ %
Andreas S., Thomas F., Alexander W., Alexander H., and Raphael W., ``ScreenshotMatcher: Taking Smartphone Photos to Capture Screenshots.'' \emph{MuC '21: Proceedings of Mensch und Computer 2021}, Ingolstadt, Germany.

\smallskip
\coloredbullet\ %
Raphael W., Leonie S., and Alexander W., ``Demonstrating GraphIT - a Platform for Working with Dependency Graphs of Learning Topics.'' \emph{23. Fachtagung Bildungstechnologien (DELFI 2025)} (\href{https://dl.gi.de/handle/20.500.12116/47137}{GI DL}).
 }

\end{document}
